Generation interstellar ships have been conceived as a solution for crewed interstellar travel aiming to bring relative numbers of humans to other star systems. Depending by the crew number generation ship’s total mass can vary wildly, constraining cruise velocity and travel time. The scope of this article is to propose a new and radical approach to human multigeneration interstellar travel in order to tackle apparently insurmountable problems that are consequent to generational interstellar travel with unconventional solutions and approaches. A comprehensive generation interstellar spaceship architecture is then proposed integrating in a unique coherent system the proposed solutions. The technologies, capabilities and approaches that are used in this proposal are all considered as currently available at TRL 3 (experimental proof of concept) and reachable in the near future, implying with that a time frame of about 50-100 years. The main driver of this proposal is to tackle four main issues that arise with long duration interstellar travel spanning centuries or (more plausibly) millennials: lack of serviceability, maintainability and reparability and degeneration of any technological system through centuries (especially electrical and computer systems), degeneration of biological and genetic variability through generations in a small populations, a reliable deceleration system for the spaceship after centuries or millennia of unusability, and, last impossibility to have a complete close-loop system for materials and resource recycling in a confined and limited environment such as a spaceship (accumulation of wastes and impossibility to reproduce new items from waste). The final purpose of this article is then to frame a realizable and robust integrated architecture from top mission objectives, requirements and constraints down to capabilities and solutions. The outcome is to prompt new breakthrough research and development in human interstellar travel and consider new technological approach to long duration space exploration which could generate further impact in research and industrial development.
